\chapter{Estado del arte}\label{cap:estado_arte}

El capítulo anterior (Capítulo~\ref{cap:marco_teorico}) ha establecido el marco formal del trabajo: la definición del problema, la familia de modelos de decisión y la formalización del problema MRTAU distribuido como un DEC-POGSMDP, junto con los algoritmos que se compararán. Durante esta exposición ya se introdujo, sobre la marcha, buena parte de la literatura relevante, este capítulo no repite las definiciones ni las descripciones algorítmicas ya descritas. Su objetivo es complementario: situar el trabajo en el estado actual de la investigación, organizar las contribuciones existentes en familias y compararlas de forma crítica frente a las particularidades del problema abordado.

Conviene además distinguir dos planos que con frecuencia se entremezclan. Uno es el de la \emph{formalización}: cómo plantear matemáticamente el problema, cuestión de la que se ocupan los modelos de decisión secuencial (Sección~\ref{sec:ea-formalizacion}). El otro es el de la \emph{resolución}: qué algoritmo computa una política sobre ese modelo (Sección~\ref{sec:ea-resolucion}). Ambos planos están relacionados, pero no deben confundirse: un mismo modelo admite múltiples métodos de resolución.

Como criterios de clasificación usaremos los ya introducidos en la Sección~\ref{sec:definicion}: el momento del cómputo (\textit{online} frente a \textit{offline}) y la arquitectura de control (centralizada frente a descentralizada). Añadiremos además como criterio adicional el tratamiento de la incertidumbre, componente esencial de todo el trabajo. Conviene acotar también el alcance: el foco se mantiene en los enfoques de planificación y coordinación para MRTA cooperativo. Quedan deliberadamente fuera las aproximaciones de aprendizaje por refuerzo multiagente, que aprenden políticas a partir de la experiencia en lugar de planificar sobre un modelo, y los enfoques de inteligencia de enjambre, por alejarse del planteamiento aplicado en este trabajo.


% ===========================================================================
\section{El problema MRTA y su evolución} \label{sec:ea-mrta}
% ===========================================================================

El problema de asignación de tareas multi-robot tiene sus raíces en la optimización combinatoria clásica. Su forma más elemental (asignar tareas a agentes minimizando un coste) es el problema de asignación resoluble en tiempo polinómico mediante el método húngaro~\cite{kuhn1955hungarian}. La incorporación de secuencias de tareas, rutas y restricciones temporales lo emparenta con el \textit{Vehicle Routing Problem} y, en presencia de ventanas de ejecución, con el VRPTW~\cite{tothvigo2014vehicle}, conexión ya señalada al describir las particularidades del escenario (Sección~\ref{subsec:mrta}).

El estudio de MRTA como disciplina propia se consolida con la taxonomía de Gerkey y Matarić~\cite{gerkey2004taxonomy}, que clasifica las variantes según los ejes ST/MT, SR/MR e IA/TA, y establece además la dureza computacional de la mayoría de ellas. Trabajos posteriores ampliaron este marco para capturar las dependencias entre tareas y agentes, en particular las \textit{cross-schedule dependencies} en las que se sitúa este problema~\cite{korsah2013taxonomy} y las restricciones temporales y de orden~\cite{nunes2017taxonomy}. Revisiones más recientes sintetizan la evolución del campo y sus principales familias de solución~\cite{khamis2015survey}. Frente a estas formulaciones, el problema abordado en este trabajo (Sección~\ref{subsec:mrta}) combina simultáneamente un conjunto de características que rara vez aparecen juntas en la literatura: tareas \emph{multi-robot} con coaliciones de tamaño $q_j$, ventanas de ejecución $[t_j^s,t_j^l]$, heterogeneidad mediante la función de compatibilidad $\textit{Cap}(\cdot,\cdot)$ y restricciones de batería con recarga. Es esta acumulación de requisitos, más que cualquiera de ellos por separado, la que hace inadecuados los enfoques diseñados para variantes más simples.

% ===========================================================================
\section{Formalización del problema: modelos de decisión} \label{sec:ea-formalizacion}
% ===========================================================================
 
El primer plano es el de la formalización: una vez delimitado el problema, ¿con qué lenguaje matemático se plantea? La respuesta predominante para la toma de decisiones secuencial bajo incertidumbre es la familia de los Procesos de Decisión de Markov presentada en la Sección~\ref{sec:modelos}. Su extensión natural al caso cooperativo con observabilidad parcial es el DEC-POMDP~\cite{bernstein2002complexity,oliehoek2016concise}, marco que captura con precisión la coordinación descentralizada con información local. La virtud de esta familia es su tratamiento de la incertidumbre y de la observabilidad parcial. Su contrapartida es la complejidad: resolver óptimamente un DEC-POMDP de horizonte finito es NEXP-completo~\cite{bernstein2002complexity}, lo que restringe la resolución exacta a instancias diminutas y obliga en la práctica a recurrir a métodos aproximados.

Para acercar estos modelos a problemas robóticos reales, una decisión fundamental para este trabajo es el uso de \emph{macro-acciones}. Una macro-acción es una acción de alto nivel y \textbf{duración variable} que encapsula internamente toda una secuencia de acciones primitivas de bajo nivel: en lugar de decidir paso a paso (avanzar, girar, etc.), el agente elige entre acciones del tipo ``desplazarse a la tarea $T_j$ y ejecutarla''. Esto tiene dos consecuencias importantes. En primer lugar, \textbf{reduce el horizonte efectivo de decisión}: el agente solo decide en los instantes en que una macro-acción termina y debe elegir la siguiente, en lugar de en cada paso elemental, lo que recorta drásticamente el número de puntos de decisión. En segundo lugar, el uso de macro-acciones obliga a introducir un modelo que sea capaz de \textbf{razonar sobre duraciones continuas}. Para ello usamos la lógica de los SMDP, que modelan los tiempos de ejecución como variables aleatorias.

El DEC-POSMDP~\cite{omidshafiei2017decentralized} y los marcos asociados de planificación con macro-acciones~\cite{amato2019modeling} explotan esta idea para integrar la descentralización y la observabilidad parcial del DEC-POMDP con las duraciones continuas del SMDP, salvando así la principal barrera que aleja a estos modelos de los problemas robóticos reales. La conexión con este trabajo es directa: las acciones locales de cada robot ($\textsc{ExecuteTask}(T_j)$, $\textsc{Recharge}$ y $\textsc{Finish}$) son precisamente macro-acciones, ya que cada una abarca desplazamiento, espera, ejecución de duración continua y estocástica en lugar de un solo movimiento atómico.

Este es, por tanto, el punto de partida del modelo definitivo propuesto en la Sección~\ref{subsec:decmrtau-decpogsmdp}. Ahora bien, las macro-acciones resuelven el problema del horizonte y de las duraciones, pero no el del \textbf{solapamiento de acciones concurrentes} entre robots, el DEC-POSMDP hereda del SMDP la incapacidad de tratar correctamente dicha concurrencia (el \emph{leak temporal} de la Sección~\ref{subsec:gsmdp}). Es esta carencia la que motiva su extensión con la lógica de eventos del GSMP, dando lugar al DEC-POGSMDP.
 
Disponer de un modelo no equivale, sin embargo, a saber resolverlo. La cuestión de qué algoritmo computa una política sobre estos modelos y, en particular, de cómo lidiar con la intratabilidad de su resolución exacta pertenece al segundo plano y es el objeto de la siguiente sección.

% ===========================================================================
\section{Métodos de resolución} \label{sec:ea-resolucion}
% ===========================================================================
 
Fijada una formalización, queda el problema de obtener la política óptima. Esta sección recorre las tres grandes familias de métodos de resolución. La primera, la programación dinámica, resuelve el modelo de forma exacta pero solo es viable en instancias diminutas; las otras dos, los métodos basados en mercado y la planificación basada en muestreo, no garantizan la optimalidad pero ofrecen soluciones en tiempos razonables de computación.
 
% --------------------------------------------------------------------------- 
\subsection{Programación dinámica exacta} \label{subsec:ea-dp}
% ---------------------------------------------------------------------------
 
El método de resolución nativo de los modelos de la sección anterior es la programación dinámica: algoritmos como la iteración de valores o la iteración de políticas calculan la política óptima resolviendo la ecuación de Bellman (Ecuación~\eqref{eq:bellman}). Para un MDP completamente observable con un espacio de estados reducido, estos métodos son exactos y eficientes. Su coste, sin embargo, crece muy rápidamente con la complejidad del modelo: la resolución exacta de un POMDP es intratable en el caso general por la dimensionalidad del espacio de creencias (Sección~\ref{subsec:pomdp}), y la de un DEC-POMDP es directamente NEXP-completa~\cite{bernstein2002complexity}. En consecuencia, la programación dinámica exacta solo resulta aplicable a instancias diminutas y funciona, en la práctica, más como referencia conceptual de optimalidad que como método utilizable a la escala del problema MRTAU. Las dos familias siguientes sacrifican esa optimalidad a cambio de escalabilidad.
 
% --------------------------------------------------------------------------- 
\subsection{Enfoques basados en mercado} \label{subsec:ea-mercado}
% ---------------------------------------------------------------------------
 
La coordinación basada en mercado traslada a los sistemas multi-robot la lógica de las subastas: las tareas se asignan a los agentes que emiten la mejor puja según una función de \emph{score}, y los conflictos se resuelven mediante negociación. Esta línea, heredera del \textit{Contract Net Protocol}, fue popularizada en robótica por sistemas como MURDOCH~\cite{gerkey2002sold} y sistematizada en la revisión de Dias et al.~\cite{dias2006market}. Su exponente descentralizado más influyente es el \emph{Consensus-Based Auction Algorithm} (CBAA) y su extensión por paquetes, el \emph{Consensus-Based Bundle Algorithm} (CBBA)~\cite{choi2009consensus}, en los que cada agente puja localmente y los conflictos se resuelven mediante una fase de consenso sobre comunicación entre vecinos. La generalización a tareas que requieren varios agentes se aborda en el \emph{Consensus-Based Grouping Algorithm} (CBGA)~\cite{hunt2014cbga}. A diferencia de la programación dinámica, estos métodos no resuelven el modelo decisión-teórico, sino que sustituyen el cálculo de la política óptima por una heurística de asignación.
 
La gran ventaja de estos métodos es que encajan a la perfección con el enfoque descentralizado: son escalables, robustos ante fallos y operan con una comunicación acotada. Su principal limitación, sin embargo, reside en que su coordinación se concentra exclusivamente en una heurística de \textit{score} de naturaleza determinista, que ni razona sobre el horizonte largo (más allá del \textit{lookahead} del paquete de CBBA) ni modela explícitamente la incertidumbre. Como consecuencia directa de ese corto horizonte de visión, el rendimiento del algoritmo depende fuertemente del diseño de dicha heurística: para garantizar un buen comportamiento a largo plazo, el \textit{score} no puede limitarse a recompensar la tarea inmediata, sino que debe incorporar componentes adicionales como fomentar la formación de coaliciones para no dejar robots esperando indefinidamente o penalizar los niveles de batería muy bajos. Esto impide adoptar directamente la propia función de recompensa del problema (Ecuación~\eqref{eq:recompensa}) como heurística y obliga a establecer prioridades e introduce constantes de ponderación, lo que convierte el ajuste de la heurística en una tarea delicada y específica del escenario. En su formulación original, además, computan una asignación estática \emph{a priori} sobre un escenario determinista. Como se detalla en la Sección~\ref{subsec:pujas}, las implementaciones de este trabajo se alejan de esos supuestos (admiten coaliciones, se ejecutan de forma \textit{online} y operan sobre un escenario estocástico mediante una aproximación determinista del \textit{score}), lo que ilustra tanto el valor como los límites de esta familia frente a un problema bajo incertidumbre.
 
% --------------------------------------------------------------------------- 
\subsection{Planificación basada en muestreo y búsqueda en árbol} \label{subsec:ea-muestreo}
% ---------------------------------------------------------------------------
 
Una tercera familia renuncia a resolver el modelo de forma cerrada y planifica mediante simulación. La búsqueda en árbol de Monte Carlo (MCTS)~\cite{browne2012survey}, y en particular su variante UCT~\cite{kocsis2006bandit} basada en el problema de la tragaperras~\cite{auer2002ucb}, construye incrementalmente un árbol de búsqueda a partir de simulaciones aleatorias. Su rasgo distintivo es que solo requiere un simulador generativo del problema, sin necesidad de la función de transición en forma cerrada, lo que la hace idónea para una dinámica definida por una cola de eventos como la del GSMDP (Sección~\ref{subsec:mrtau-gsmdp}). Su carácter \textit{anytime} la alinea, además, con la necesidad de cómputo acotado de los métodos \textit{online} (Sección~\ref{subsec:online-offline}).
 
La adaptación de MCTS al contexto multi-robot descentralizado la proporciona Dec-MCTS~\cite{best2019decmcts}, en el que cada robot mantiene su propio árbol y comunica a sus vecinos una distribución de probabilidad sobre sus planes, empleando una variante descontada de UCT (D-UCT) para adaptarse a los planes no estacionarios ajenos~\cite{garivier2011ducb}. Esta representación probabilística de las intenciones de los compañeros constituye la forma de coordinación más rica de las analizadas y es la base de la implementación descrita en la Sección~\ref{subsec:dec-mcts}. Su precio es un coste de cómputo elevado y una implementación delicada, lo que justifica su comparación empírica frente a las familias anteriores. El coste de cómputo introducido queda plenamente justificado por la naturaleza del problema. Debido a su carácter \textit{online}, el tiempo requerido para desplazarse, recargar o ejecutar una tarea puede utilizarse como tiempo de cómputo para enriquecer una próxima decisión en lugar de ser tiempo ``perdido'' como en las aproximaciones \textit{on-request}.

% ===========================================================================
\section{Asignación de tareas bajo incertidumbre} \label{sec:ea-incertidumbre}
% ===========================================================================

La inmensa mayoría de la literatura clásica de MRTA, incluidas las taxonomías~\cite{gerkey2004taxonomy,korsah2013taxonomy}, asume un escenario \emph{determinista}. El tratamiento integrado de la incertidumbre está comparativamente menos explorado y constituye un área activa. El trabajo de Choudhury et al.~\cite{choudhury2020dynamic} sobre asignación dinámica de tareas bajo incertidumbre y restricciones temporales es, por su planteamiento, uno de los más cercanos en espíritu al abordado aquí.

Las herramientas formales para tratar esta incertidumbre se han presentado en el Capítulo~\ref{cap:marco_teorico}: los modelos de decisión bajo observabilidad parcial y, de forma central, los procesos semi-Markovianos generalizados~\cite{glynn1989gsmp,younes2004gsmdp}, que modela el solapamiento de acciones de duración continua mediante eventos con relojes ocultos. Lo que la literatura rara vez ofrece es la \emph{combinación} de estas herramientas con el resto de exigencias del problema: descentralización plena, formación de coaliciones, ventanas de ejecución y restricciones de batería, todo ello de manera simultánea. En particular, el problema de la fuga de información temporal (Sección~\ref{subsec:gsmdp}), que surge al modelar duraciones continuas concurrentes bajo observabilidad total, no suele plantearse de forma explícita en el contexto de MRTA, pese a ser esencial para la corrección del modelo.


% ===========================================================================
\section{Posicionamiento y aportaciones} \label{sec:ea-posicionamiento}
% ===========================================================================

En el plano de la formalización, los modelos de decisión (Sección~\ref{sec:ea-formalizacion}) describen el problema pero su resolución exacta mediante programación dinámica (Sección~\ref{subsec:ea-dp}) es generalmente intratable. En el plano de la resolución aproximada, los métodos basados en mercado (Sección~\ref{subsec:ea-mercado}) son escalables y descentralizados, pero condensan su coordinación en heurísticas deterministas con escasa anticipación, mientras que la planificación basada en muestreo (Sección~\ref{subsec:ea-muestreo}) ofrece un punto intermedio \textit{anytime} y descentralizable, a costa de un cómputo elevado.
 
Sobre este análisis de pueden definir las principales aportaciones de este trabajo. En primer lugar, la \textbf{formalización del problema MRTAU distribuido como un DEC-POGSMDP} (Sección~\ref{subsec:decmrtau-decpogsmdp}): un modelo que, hasta donde se ha podido constatar, no aparece en la literatura bajo este nombre y que se introduce como la progresión natural que resulta de dotar a un DEC-POSMDP de la lógica de eventos del GSMP, garantizando así el tratamiento correcto de las duraciones continuas concurrentes y la ausencia de \emph{leak temporal}. En segundo lugar, la \textbf{definición unificada de los algoritmos como políticas de comunicación} que resuelven ese mismo DEC-POGSMDP con grados crecientes de coordinación (Sección~\ref{sec:algoritmos}), lo que permite comparar en igualdad de condiciones desde los \textit{baselines} sin comunicación hasta la planificación distribuida de Dec-MCTS. En tercer lugar, y como hilo conductor de toda la memoria, la \textbf{transición desde el modelo centralizado de partida} (el MCTS centralizado de~\red{CITA AL TUTOR}) \textbf{hacia una resolución puramente distribuida}: una reformulación que, como se argumentó en la Sección~\ref{sec:formalizacion}, no es un mero detalle de implementación, sino un cambio en la clase de complejidad del problema que obliga a sustituir la información global por mecanismos de coordinación basados en comunicación.